% Generated draft: requires mathematical and source review.
\documentclass[11pt]{article}
\usepackage[margin=25mm]{geometry}
\usepackage{fontspec}
\setmainfont{Latin Modern Roman}
\usepackage{amsmath,amssymb,mathtools,unicode-math}
\setmathfont{Latin Modern Math}
\usepackage{enumitem,needspace,xurl,hyperref}
\setlength{\parindent}{0pt}
\setlength{\parskip}{5pt}
\setlength{\emergencystretch}{3em}
\newcommand{\N}{\mathbb{N}}
\newcommand{\Z}{\mathbb{Z}}
\newcommand{\Q}{\mathbb{Q}}
\newcommand{\R}{\mathbb{R}}
\newcommand{\C}{\mathbb{C}}
\newcommand{\F}{\mathbb{F}}
\newcommand{\A}{\mathbb{A}}
\newcommand{\PP}{\mathbb{P}}
\newcommand{\E}{\mathbb{E}}
\newcommand{\eps}{\varepsilon}
\newcommand{\dd}{\,\mathrm{d}}
\begin{document}
\section*{UnsolvedMath ID: 30001514}
\textbf{Source problem code:} catalogue entry
\subsection{Definitions and mathematical statement}
The supplied problem asks for conditions on the domain and boundary data of the affine Plateau problem that ensure boundary attainment of the gradient, described in the record as equivalent to appropriate boundary behavior of the Gauss map. The record does not specify the formulation of the affine Plateau problem, its admissible solutions, or the prescribed boundary data, so a fully quantified mathematical statement cannot be recovered from it.

For clarity, in a conditional graph interpretation, let \(\Omega\subset\R^n\) be an open set, let \(u\in C^1(\Omega)\), and let \(G:\partial\Omega\to\R^n\) denote a prescribed boundary gradient field. Pointwise boundary gradient attainment would mean
\[
\forall x_0\in\partial\Omega,\qquad
\lim_{\substack{x\to x_0\\x\in\Omega}}\nabla u(x)=G(x_0).
\]
These objects and this interpretation are explanatory only: the supplied record does not establish that its boundary data are a vector field \(G\), that its solutions have this regularity, or that pointwise attainment at every boundary point is intended.

\emph{Source question.} Determine hypotheses on the domain and boundary data that guarantee the intended boundary gradient attainment for solutions of the affine Plateau problem. The corresponding Gauss-map boundary condition and the claimed equivalence require specification of the geometric formulation and the meaning of boundary attainment. No particular dimension, domain geometry, regularity class, solution concept, or mode of convergence is supplied; none is imposed here as a hypothesis of the source problem.
\subsection{Short English statement}
Find conditions on the domain and boundary data that force solutions of the affine Plateau problem to attain the prescribed gradient behavior at the boundary. The source also describes this through the Gauss map, but leaves the precise formulation and boundary convergence unspecified. This is a statement-only draft requiring human mathematical review.
\subsection{Sources}
\begin{itemize}
\item Local registry: \texttt{data/unsolvedmath/problems.json} and the supplementary registry \texttt{catalog/problems.json}; supplied record with permanent UnsolvedMath ID \(30001514\), source code \texttt{catalogue entry}, and title \emph{Boundary Gradient Attainment in the Affine Plateau Problem}. These are registry identifiers, not independently inspected source locations.
\item Supplied report citation and literature-source link: \emph{Calculus of Variations (2011)}, Oberwolfach Reports, Report 31, pp.~1827--1884, DOI \texttt{10.4171/owr/2010/31}, \url{https://doi.org/10.4171/owr/2010/31}. The collection attributes the problem to workshop open-problem discussions documented in Oberwolfach Reports; no individual proposer is supplied. The proposed year is recorded as 2010.
\item Supplied literature assessment, checked \texttt{2026-08-21}: the registry labels the problem open, reports that the report explicitly identifies boundary gradient or Gauss-map behavior as open, and states that the literature checked did not identify a result settling the question in this generality. The listed checked source is the report linked above. This is supplied dated curation derived from the OWR collection, not an independently verified assessment; time-sensitive status and cited claims require independent verification.
\item The registry labels the statement \texttt{exact}; its verification note says that the statement is internally readable and retained byte-for-byte. This supplied label does not establish correspondence with the report or mathematical completeness.
\end{itemize}
\end{document}
